% !TeX root = ../tfg.tex
% !TeX encoding = utf8

\newcommand{\bestsicacivae}{mlp_vae_sicap_c1_s42_lr0.0001}
\newcommand{\bestsicaciddpm}{mlp_ddpm_sicap_c1_s42_lr0.0001_t10}
\newcommand{\bestsicacxiivae}{mlp_vae_sicap_c12_s42_lr0.0005}
\newcommand{\bestsicacxiiddpm}{mlp_ddpm_sicap_c12_s42_lr0.001_t10}

\chapter{Resultados y discusión}
\label{ch:resultados}

Este capítulo presenta los resultados experimentales siguiendo la progresión de
dificultad descrita en el capítulo~\ref{ch:experimental}: comienza con los conjuntos
sintéticos como validación cualitativa del comportamiento esperado de los modelos,
continúa con los conjuntos de imágenes MNIST y PathMNIST como puente entre el dominio
controlado y el dominio real, y se concluye con SICAP como caso de aplicación de mayor
relevancia clínica. Para cada conjunto se analizan primero las representaciones
aprendidas (espacio latente para el VAE, reconstrucciones para el DDPM) y después las
distribuciones de \textit{scores} y las curvas ROC que cuantifican la capacidad de
detección OOD. El capítulo cierra con un análisis comparativo entre los métodos
generativos propuestos y las líneas base de distancia directa en el espacio de
características, seguido de una discusión de las limitaciones del trabajo.

\section{Resultados de los experimentos}
\label{sec:res-experimentos}

\subsection{Conjuntos sintéticos}
\label{subsec:res-sinteticos}

Los conjuntos sintéticos proporcionan un entorno en el que el comportamiento
esperado de los modelos puede verificarse visualmente: si un modelo generativo ha
aprendido correctamente la distribución ID, sus reconstrucciones de muestras OOD
deben diferir de las originales. Las
figuras~\ref{fig:recon_vae_blobs} y~\ref{fig:recon_vae_moons} muestran los espacios
de entrada y reconstruido del VAE para Blobs y Moons, respectivamente. En ambos casos
se observa que las muestras ID se reconstruyen fielmente, mientras que las muestras
OOD quedan desplazadas hacia la distribución de entrenamiento: el modelo actúa como
un atractor hacia la variedad aprendida. Este efecto es más pronunciado en Blobs,
donde la separación lineal hace que las reconstrucciones OOD se desplacen claramente
hacia el clúster ID, mientras que en Moons la reconstrucción OOD queda en la zona de
transición debido a la mayor proximidad geométrica entre distribuciones. El comportamiento del DDPM resulta análogo, por lo que prescindiremos de sus figuras.

\begin{figure}[htbp]
    \centering
    \includegraphics[width=\textwidth]{results-evaluation-mlp_vae_moons_s42_lr0.001-plots-reconstructions.pdf}
    \caption{Espacio de entrada (izquierda) y espacio reconstruido por el VAE
             (derecha) en el conjunto \textit{Moons}. Las reconstrucciones OOD
             se desplazan hacia la variedad de entrenamiento, generando un error
             de reconstrucción sistemáticamente mayor que el de las muestras ID.}
    \label{fig:recon_vae_moons}
\end{figure}

\begin{figure}[htbp]
    \centering
    \includegraphics[width=\textwidth]{results-evaluation-mlp_vae_blobs_s42_lr0.001-plots-reconstructions.pdf}
    \caption{Espacio de entrada y espacio reconstruido del VAE en \textit{Blobs}.
             La separación lineal entre clústeres facilita la detección OOD incluso
             con un modelo de baja capacidad ($L{=}2$), lo que sirve como caso de
             referencia para comparar con \textit{Moons}.}
    \label{fig:recon_vae_blobs}
\end{figure}

\subsection{VAE: MNIST y PathMNIST}
\label{subsec:res-vae-mnist-path}

Las figuras~\ref{fig:vae_recon_mnist} y~\ref{fig:vae_recon_pathmnist} muestran
muestras originales y reconstruidas por el VAE para MNIST/SVHN y PathMNIST
respectivamente. Para MNIST, las reconstrucciones de dígitos ID son limpias y nítidas,
mientras que las imágenes OOD~(SVHN) se reconstruyen con mayor borrosidad y pérdida
de detalle, inclusive presentando algunos \emph{features} generales de MNIST.

Para PathMNIST (figura~\ref{fig:vae_recon_pathmnist}), las diferencias entre las
reconstrucciones ID y OOD son más sutiles: ambas presentan una calidad visual similar,
lo que anticipa que el \textit{score} de reconstrucción MSE por sí solo puede ser
insuficiente para discriminar ambas distribuciones. La mayor dificultad de este dominio
respecto a MNIST se origina en que los tejidos ID y OOD (mucosa normal, estroma y
epitelio tumoral) comparten patrones de textura a escalas intermedias que son
difícilmente capturables con un espacio latente de dimensión acotada.

Las proyecciones PCA y UMAP del espacio latente
(figuras~\ref{fig:vae_embedding_mnist} y~\ref{fig:vae_embedding_pathmnist}) ilustran
la separabilidad que el modelo consigue en la representación interna. Para MNIST, la
separación entre la nube ID y los puntos OOD es clara incluso en la proyección lineal
PCA, lo que implica que los \textit{scores} basados en distancia latente (k-NN y
Mahalanobis) disponen de una representación discriminativa en la que apoyarse. Para
PathMNIST, las nubes ID y OOD presentan mayor solapamiento en las proyecciones
latentes, lo que anuncia valores de AUROC más moderados para los \textit{scores}
basados en la geometría del espacio latente. Los resultados completos de las
distribuciones de \textit{scores} y las curvas ROC para MNIST y PathMNIST se
recogen en el apéndice~\ref{ch:apendice-figuras}.

\begin{figure}[htbp]
    \centering
    \includegraphics[width=\textwidth]{results-evaluation-mlp_vae_mnist_s42_lr0.0001-plots-reconstructions_compact.pdf}
    \caption{Muestras originales y reconstruidas por el VAE en MNIST/SVHN.
             Las reconstrucciones OOD presentan mayor borrosidad que las ID,
             respaldando el uso del MSE como \textit{score} de anomalía.}
    \label{fig:vae_recon_mnist}
\end{figure}

\begin{figure}[htbp]
    \centering
    \includegraphics[width=\textwidth]{results-evaluation-mlp_vae_pathmnist_c1_s42_lr0.0001-plots-reconstructions_compact.pdf}
    \caption{Muestras originales y reconstruidas por el VAE en PathMNIST~(tejido
             normal ID frente a tumoral OOD). Las diferencias de textura son
             más sutiles que en MNIST, reflejando la mayor dificultad del dominio
             histopatológico.}
    \label{fig:vae_recon_pathmnist}
\end{figure}

\begin{figure}[htbp]
    \centering
    \includegraphics[width=\textwidth]{results-evaluation-mlp_vae_mnist_s42_lr0.0001-plots-embedding_vs_ood.pdf}
    \caption{Espacio latente del VAE en MNIST (proyecciones PCA y UMAP) con
             datos ID y OOD superpuestos. La separación latente es condición
             necesaria para que los \textit{scores} $k$-NN y Mahalanobis
             latentes sean efectivos.}
    \label{fig:vae_embedding_mnist}
\end{figure}

\begin{figure}[htbp]
    \centering
    \includegraphics[width=\textwidth]{results-evaluation-mlp_vae_pathmnist_c1_s42_lr0.0001-plots-embedding_vs_ood.pdf}
    \caption{Espacio latente del VAE en PathMNIST (proyecciones PCA y UMAP).
             La mayor superposición entre nubes ID y OOD anticipa valores de AUROC
             más moderados para los \textit{scores} basados en distancia latente.}
    \label{fig:vae_embedding_pathmnist}
\end{figure}

\subsection{DDPM: MNIST y PathMNIST}
\label{subsec:res-ddpm-mnist-path}

Las figuras~\ref{fig:ddpm_tgrid_mnist} y~\ref{fig:ddpm_tgrid_pathmnist} muestran la
reconstrucción iterativa del DDPM a distintos niveles de corrupción~$t$ para una muestra
ID y una OOD de cada conjunto: se añade ruido hasta el instante~$t$ y se ejecuta la
cadena inversa paso a paso hasta~$t{=}0$. Para MNIST, incluso a valores de~$t$ moderados
la muestra ID se reconstruye fielmente, mientras que la muestra SVHN (OOD) acumula un
error creciente a medida que aumenta la corrupción: el proceso de \textit{denoising}
no puede recuperar los rasgos fotorrealistas del dígito SVHN porque el modelo nunca
los ha visto durante el entrenamiento y los reemplaza por patrones manuscritos. Para
PathMNIST (figura~\ref{fig:ddpm_tgrid_pathmnist}), la diferencia visual entre las
reconstrucciones ID y OOD es más sutil, lo que es consistente con la mayor proximidad
distribucional de los tejidos implicados.

Las figuras~\ref{fig:ddpm_traj_mnist} y~\ref{fig:ddpm_traj_pathmnist} muestran las
trayectorias de \textit{denoising} inverso para ambos conjuntos, partiendo del mismo
nivel de ruido inicial. La divergencia entre la trayectoria ID y la trayectoria OOD
ilustra cualitativamente la diferencia de \textit{score} asignado: la muestra ID
converge hacia una reconstrucción fiel de la imagen original, mientras que la muestra
OOD converge hacia una imagen que pertenece a la distribución ID aprendida pero no
coincide con la muestra original. Este comportamiento es la manifestación visual del
mecanismo de puntuación: el modelo no puede \textit{denoisear} la OOD hacia sí misma
porque su distribución no está en el soporte aprendido, y el error entre el resultado
y la entrada original actúa como señal de anomalía. Los resultados cuantitativos para
estos conjuntos se recogen en el apéndice~\ref{ch:apendice-figuras}.

\begin{figure}[htbp]
    \centering
    \includegraphics[width=\textwidth]{results-evaluation-mlp_ddpm_mnist_s42_lr0.0001_t10-plots-ddpm_timestep_grid_compact.pdf}
    \caption{Reconstrucción iterativa del DDPM a distintos niveles de
             corrupción~$t$ para una muestra MNIST~(ID) y una SVHN~(OOD). El
             mayor error de reconstrucción en datos OOD es visible incluso a~$t$
             bajos.}
    \label{fig:ddpm_tgrid_mnist}
\end{figure}

\begin{figure}[htbp]
    \centering
    \includegraphics[width=\textwidth]{results-evaluation-mlp_ddpm_pathmnist_c1_s42_lr0.0001_t10-plots-ddpm_timestep_grid_compact.pdf}
    \caption{Reconstrucción iterativa del DDPM a distintos niveles de
             corrupción~$t$ para una muestra PathMNIST normal~(ID) y tumoral~(OOD).
             Las diferencias son más sutiles que en MNIST, reflejando la dificultad
             del dominio histopatológico.}
    \label{fig:ddpm_tgrid_pathmnist}
\end{figure}

\begin{figure}[htbp]
    \centering
    \includegraphics[width=\textwidth]{results-evaluation-mlp_ddpm_mnist_s42_lr0.0001_t10-plots-ddpm_denoising_trajectory_compact.pdf}
    \caption{Trayectoria de \textit{denoising} inverso del DDPM para una muestra
             MNIST~(ID) y una SVHN~(OOD) desde el mismo nivel de ruido inicial. La
             divergencia entre trayectorias ilustra cualitativamente la diferencia
             de \textit{score} OOD asignado.}
    \label{fig:ddpm_traj_mnist}
\end{figure}

\begin{figure}[htbp]
    \centering
    \includegraphics[width=\textwidth]{results-evaluation-mlp_ddpm_pathmnist_c1_s42_lr0.0001_t10-plots-ddpm_denoising_trajectory_compact.pdf}
    \caption{Trayectoria de \textit{denoising} inverso del DDPM en PathMNIST
             para tejido normal~(ID) y tumoral~(OOD). La evolución de ambas
             muestras bajo el proceso de \textit{denoising} refleja la diferencia
             en su compatibilidad con la distribución aprendida.}
    \label{fig:ddpm_traj_pathmnist}
\end{figure}

\subsection{SICAP: resultados del mejor experimento por conjunto}
\label{subsec:res-sicap}

SICAP es el conjunto de mayor dificultad y relevancia clínica de los estudiados.
Los resultados que se presentan a continuación corresponden a la configuración con
mayor AUROC de validación para cada combinación de modelo y conjunto (C1 o C12),
tal y como se recoge en las macros definidas al inicio del capítulo. La comparación
sistemática entre todas las configuraciones se presenta en la
sección~\ref{sec:res-comparativo}.

\paragraph{VAE en SICAP.}
Las figuras~\ref{fig:sicapc1_vae_recon}--\ref{fig:sicapc1_vae_lmah} muestran los
cuatro \textit{scores} del VAE para SICAP-C1: error de reconstrucción~(MSE),
ELBO, $k$-NN latente y distancia de Mahalanobis latente. Para cada \textit{score}
se representa la distribución de puntuaciones sobre el conjunto \texttt{id\_eval}
y sobre \texttt{ood\_eval}, junto con la curva ROC correspondiente. El grado de
separación entre histogramas y el área bajo la curva ROC revelan que los
\textit{scores} latentes (k-NN y Mahalanobis) tienden a ser más discriminativos
que los \textit{scores} en el espacio de píxeles (MSE y ELBO), lo que refleja que
el espacio latente del VAE captura la estructura distribucional de la ID de forma
más compacta que el espacio de entrada. No obstante, el considerable solapamiento
entre las distribuciones de puntuaciones ID y OOD es coherente con el solapamiento
observado en el espacio de características UNI
(figura~\ref{fig:input_space_sicap}): cuando las distribuciones ID y OOD son
similares en el espacio de entrada, ninguna transformación determinista (como la
codificación latente) puede separarlas perfectamente.

Las figuras~\ref{fig:sicapc12_vae_recon}--\ref{fig:sicapc12_vae_lmah} muestran los
resultados equivalentes para SICAP-C12. El ensanchamiento de la distribución ID al
incorporar Gleason~2 al entrenamiento tiende a reducir la separación respecto a OOD
(Gleason~3--4), dado que la diferencia estadística entre distribuciones disminuye.
Sin embargo, en algunos \textit{scores}, la mayor cobertura ID de SICAP-C12 permite
al modelo aprender una representación latente más rica de la normalidad, lo que puede
mejorar la capacidad discriminativa frente a anomalías de alto grado.

\paragraph{DDPM en SICAP.}
Las figuras~\ref{fig:sicapc1_ddpm_single}--\ref{fig:sicapc1_ddpm_rknn} presentan los
cinco \textit{scores} del DDPM para SICAP-C1. El \textit{score single-step} MSE
(figura~\ref{fig:sicapc1_ddpm_single}) mide el error de predicción del ruido a un
único instante fijo~$t^*$: es el \textit{score} más simple computacionalmente y
su discriminabilidad depende de la elección del instante de evaluación. Los
\textit{scores multi-step} MSE y coseno
(figuras~\ref{fig:sicapc1_ddpm_multimse}--\ref{fig:sicapc1_ddpm_multicos}) promedian
el error normalizado por $z$-\textit{score} sobre~$T{=}10$ instantes de tiempo
uniformemente espaciados entre el $10\,\%$ y el $90\,\%$ del rango $[1, T_{\max}]$,
reduciendo la sensibilidad al valor de~$t^*$ y proporcionando una estimación más
robusta de la anomalía. Los \textit{scores} basados en el residuo de predicción de
ruido (Mahalanobis y k-NN, figuras~\ref{fig:sicapc1_ddpm_rmah}--\ref{fig:sicapc1_ddpm_rknn})
evalúan la distancia del residuo $(\hat{\varepsilon} - \varepsilon)$ respecto a la
distribución de residuos del entrenamiento, y pueden ser más discriminativos cuando
el DDPM comete errores sistemáticamente distintos sobre muestras OOD. Las
figuras~\ref{fig:sicapc12_ddpm_single}--\ref{fig:sicapc12_ddpm_rknn} muestran los
resultados equivalentes para SICAP-C12.

% ── SICAP-C1 · VAE ───────────────────────────────────────────────────────────

\begin{figure}[htbp]
    \centering
    \begin{minipage}[t]{0.48\textwidth}
        \includegraphics[width=\textwidth]{results-evaluation-\bestsicacivae-plots-ood_scores_recon.pdf}
        \caption{\textit{Score} MSE de reconstrucción del mejor VAE en SICAP-C1
                 (ID: Gleason~1; OOD: Gleason~2--4).}
        \label{fig:sicapc1_vae_recon}
    \end{minipage}
    \hfill
    \begin{minipage}[t]{0.48\textwidth}
        \includegraphics[width=\textwidth]{results-evaluation-\bestsicacivae-plots-ood_scores_elbo.pdf}
        \caption{\textit{Score} ELBO del mejor VAE en SICAP-C1.}
        \label{fig:sicapc1_vae_elbo}
    \end{minipage}
\end{figure}

\begin{figure}[htbp]
    \centering
    \begin{minipage}[t]{0.48\textwidth}
        \includegraphics[width=\textwidth]{results-evaluation-\bestsicacivae-plots-ood_scores_latent_knn.pdf}
        \caption{\textit{Score} $k$-NN latente del mejor VAE en SICAP-C1.}
        \label{fig:sicapc1_vae_lknn}
    \end{minipage}
    \hfill
    \begin{minipage}[t]{0.48\textwidth}
        \includegraphics[width=\textwidth]{results-evaluation-\bestsicacivae-plots-ood_scores_latent_mah.pdf}
        \caption{\textit{Score} de Mahalanobis latente del mejor VAE en SICAP-C1.}
        \label{fig:sicapc1_vae_lmah}
    \end{minipage}
\end{figure}

% ── SICAP-C1 · DDPM ──────────────────────────────────────────────────────────

\begin{figure}[htbp]
    \centering
    \begin{minipage}[t]{0.48\textwidth}
        \includegraphics[width=\textwidth]{results-evaluation-\bestsicaciddpm-plots-ood_scores_noise_single.pdf}
        \caption{\textit{Score} \textit{single-step} MSE del mejor DDPM en SICAP-C1.}
        \label{fig:sicapc1_ddpm_single}
    \end{minipage}
    \hfill
    \begin{minipage}[t]{0.48\textwidth}
        \includegraphics[width=\textwidth]{results-evaluation-\bestsicaciddpm-plots-ood_scores_noise_multi_mse.pdf}
        \caption{\textit{Score} \textit{multi-step} MSE ($z$-\textit{score}) del
                 mejor DDPM en SICAP-C1.}
        \label{fig:sicapc1_ddpm_multimse}
    \end{minipage}
\end{figure}

\begin{figure}[htbp]
    \centering
    \begin{minipage}[t]{0.48\textwidth}
        \includegraphics[width=\textwidth]{results-evaluation-\bestsicaciddpm-plots-ood_scores_noise_multi_cosine.pdf}
        \caption{\textit{Score} \textit{multi-step} coseno ($z$-\textit{score})
                 del mejor DDPM en SICAP-C1.}
        \label{fig:sicapc1_ddpm_multicos}
    \end{minipage}
    \hfill
    \begin{minipage}[t]{0.48\textwidth}
        \includegraphics[width=\textwidth]{results-evaluation-\bestsicaciddpm-plots-ood_scores_residual_mah.pdf}
        \caption{\textit{Score} de Mahalanobis sobre el residuo de predicción de
                 ruido del mejor DDPM en SICAP-C1.}
        \label{fig:sicapc1_ddpm_rmah}
    \end{minipage}
\end{figure}

\begin{figure}[htbp]
    \centering
    \includegraphics[width=\textwidth]{results-evaluation-\bestsicaciddpm-plots-ood_scores_residual_knn.pdf}
    \caption{Distribución del \textit{score} $k$-NN sobre el residuo de predicción
             de ruido y curva ROC del mejor experimento DDPM en SICAP-C1.}
    \label{fig:sicapc1_ddpm_rknn}
\end{figure}

% ── SICAP-C12 · VAE ──────────────────────────────────────────────────────────

\begin{figure}[htbp]
    \centering
    \begin{minipage}[t]{0.48\textwidth}
        \includegraphics[width=\textwidth]{results-evaluation-\bestsicacxiivae-plots-ood_scores_recon.pdf}
        \caption{\textit{Score} MSE de reconstrucción del mejor VAE en SICAP-C12
                 (ID: Gleason~1--2; OOD: Gleason~3--4).}
        \label{fig:sicapc12_vae_recon}
    \end{minipage}
    \hfill
    \begin{minipage}[t]{0.48\textwidth}
        \includegraphics[width=\textwidth]{results-evaluation-\bestsicacxiivae-plots-ood_scores_elbo.pdf}
        \caption{\textit{Score} ELBO del mejor VAE en SICAP-C12.}
        \label{fig:sicapc12_vae_elbo}
    \end{minipage}
\end{figure}

\begin{figure}[htbp]
    \centering
    \begin{minipage}[t]{0.48\textwidth}
        \includegraphics[width=\textwidth]{results-evaluation-\bestsicacxiivae-plots-ood_scores_latent_knn.pdf}
        \caption{\textit{Score} $k$-NN latente del mejor VAE en SICAP-C12.}
        \label{fig:sicapc12_vae_lknn}
    \end{minipage}
    \hfill
    \begin{minipage}[t]{0.48\textwidth}
        \includegraphics[width=\textwidth]{results-evaluation-\bestsicacxiivae-plots-ood_scores_latent_mah.pdf}
        \caption{\textit{Score} de Mahalanobis latente del mejor VAE en SICAP-C12.}
        \label{fig:sicapc12_vae_lmah}
    \end{minipage}
\end{figure}

% ── SICAP-C12 · DDPM ─────────────────────────────────────────────────────────

\begin{figure}[htbp]
    \centering
    \begin{minipage}[t]{0.48\textwidth}
        \includegraphics[width=\textwidth]{results-evaluation-\bestsicacxiiddpm-plots-ood_scores_noise_single.pdf}
        \caption{\textit{Score} \textit{single-step} MSE del mejor DDPM en SICAP-C12.}
        \label{fig:sicapc12_ddpm_single}
    \end{minipage}
    \hfill
    \begin{minipage}[t]{0.48\textwidth}
        \includegraphics[width=\textwidth]{results-evaluation-\bestsicacxiiddpm-plots-ood_scores_noise_multi_mse.pdf}
        \caption{\textit{Score} \textit{multi-step} MSE ($z$-\textit{score}) del
                 mejor DDPM en SICAP-C12.}
        \label{fig:sicapc12_ddpm_multimse}
    \end{minipage}
\end{figure}

\begin{figure}[htbp]
    \centering
    \begin{minipage}[t]{0.48\textwidth}
        \includegraphics[width=\textwidth]{results-evaluation-\bestsicacxiiddpm-plots-ood_scores_noise_multi_cosine.pdf}
        \caption{\textit{Score} \textit{multi-step} coseno ($z$-\textit{score})
                 del mejor DDPM en SICAP-C12.}
        \label{fig:sicapc12_ddpm_multicos}
    \end{minipage}
    \hfill
    \begin{minipage}[t]{0.48\textwidth}
        \includegraphics[width=\textwidth]{results-evaluation-\bestsicacxiiddpm-plots-ood_scores_residual_mah.pdf}
        \caption{\textit{Score} de Mahalanobis sobre el residuo de predicción de
                 ruido del mejor DDPM en SICAP-C12.}
        \label{fig:sicapc12_ddpm_rmah}
    \end{minipage}
\end{figure}

\begin{figure}[htbp]
    \centering
    \includegraphics[width=\textwidth]{results-evaluation-\bestsicacxiiddpm-plots-ood_scores_residual_knn.pdf}
    \caption{Distribución del \textit{score} $k$-NN sobre el residuo de predicción
             de ruido y curva ROC del mejor experimento DDPM en SICAP-C12.}
    \label{fig:sicapc12_ddpm_rknn}
\end{figure}

% ─────────────────────────────────────────────────────────────────────────────
\section{Análisis comparativo entre métodos}
\label{sec:res-comparativo}
% ─────────────────────────────────────────────────────────────────────────────

La tabla~\ref{tab:sicap_comparison} y la figura~\ref{fig:sicap_methods_comparison}
resumen el AUROC del mejor experimento de cada familia de métodos en SICAP-C1 y
SICAP-C12: los métodos de distancia directa en el espacio de características UNI
(k-NN y Mahalanobis sobre las features crudas), el VAE y el DDPM. Esta comparación
es especialmente informativa porque todos los métodos operan sobre la misma
representación de entrada (características UNI de dimensión~$1\,024$), lo que elimina
el efecto de la representación y permite aislar la contribución del modelo generativo.

Los métodos de distancia directa representan la línea base sin entrenamiento generativo:
simplemente construyen una referencia de normalidad sobre las características de
entrenamiento y evalúan la distancia de cada muestra test respecto a esa referencia.
Su sencillez conceptual y computacional los convierte en un estándar comparativo natural;
si los modelos generativos no los superan consistentemente, su coste adicional de
entrenamiento no estaría justificado. La tabla~\ref{tab:sicap_comparison} permite
contrastar directamente si el aprendizaje de una distribución generativa en el espacio
de características UNI aporta una ganancia de discriminabilidad respecto a operar
directamente con distancias en ese espacio.

Los resultados muestran que la ventaja relativa de los modelos generativos depende
del conjunto y del modo de puntuación. Para SICAP-C1, los \textit{scores} basados
en la geometría del espacio latente del VAE (k-NN y Mahalanobis latentes) tienden a
competir con los métodos de distancia directa, dado que ambos operan en espacios de
baja dimensión con estructura gaussiana implícita. Para el DDPM, los modos
\textit{multi-step} (que integran información de múltiples instantes temporales)
aportan una perspectiva complementaria que puede ser más o menos discriminativa
dependiendo de si el modelo ha logrado capturar la estructura de la distribución ID
en el espacio de ruido predicho. La figura~\ref{fig:sicap_methods_comparison}, que
muestra el AUROC medio con barras de error sobre distintas semillas, permite evaluar
la estabilidad de cada método respecto a la inicialización aleatoria.

\begin{table}[htbp]
    \centering
    \caption{Comparación entre métodos en SICAP: distancias de referencia
             (KNN, Mahalanobis) frente a modelos generativos (VAE, DDPM).
             Se muestra la configuración con mayor AUROC de cada familia;
             se resalta en negrita el mejor resultado global por conjunto.
             La tabla completa con todos los hiperparámetros figura en el
             apéndice~\ref{ch:apendice-figuras} (tabla~\ref{tab:sicap_results_full}).}
    \label{tab:sicap_comparison}
    \input{tables/table_comparison_sicap.tex}
\end{table}

\begin{figure}[htbp]
    \centering
    \includegraphics[width=\textwidth]{results-summary-sicap_methods_comparison.pdf}
    \caption{AUROC medio por método en SICAP-C1 y SICAP-C12. Las barras de
             error corresponden a la desviación estándar sobre distintas semillas
             de inicialización. Los métodos de distancia directa (KNN, Mahalanobis)
             actúan como línea base sin entrenamiento generativo.}
    \label{fig:sicap_methods_comparison}
\end{figure}

\paragraph{Ablación sobre el número de pasos de puntuación $T$.}~\\

Una de las variables de diseño más relevantes para el DDPM es el número de
instantes~$T$ sobre los que se promedia el error normalizado en los modos
\textit{multi-step}. Incrementar~$T$ proporciona una estimación más robusta de la
anomalía al integrar información de un mayor rango del proceso de difusión, pero
aumenta el coste computacional de la evaluación en un factor~$T$ respecto al modo
\textit{single-step}. La tabla~\ref{tab:ablation_t} y la
figura~\ref{fig:ablation_t_auroc} muestran el AUROC en función de~$T$ para los tres
\textit{scores} DDPM ($S_{\text{single}}$, $S_{\text{multi-MSE}}$, $S_{\text{multi-cos}}$)
en SICAP-C1 y SICAP-C12.

Los resultados revelan que la ganancia de AUROC al pasar de~$T{=}1$ a~$T{=}10$ es
sustancial en los modos \textit{multi-step}, mientras que incrementos adicionales
($T{>}10$) producen mejoras marginales que difícilmente justifican el coste
computacional añadido. Esta saturación sugiere que los primeros~$T{=}10$ pasos ya
capturan la mayor parte de la información discriminativa disponible en el proceso de
difusión. El modo \textit{single-step} es insensible a~$T$ por definición y sirve
como referencia de eficiencia computacional mínima.

\begin{figure}[htbp]
    \centering
    \includegraphics[width=\textwidth]{results-summary-ablation_t_auroc_sicap.pdf}
    \caption{AUROC en función del número de pasos de puntuación~$T$ del DDPM
             para los tres \textit{scores} (\textit{single-step} MSE,
             \textit{multi-step} MSE y \textit{multi-step} coseno) en SICAP-C1
             y SICAP-C12. Las barras de error muestran la variabilidad entre
             semillas.}
    \label{fig:ablation_t_auroc}
\end{figure}

\begin{table}[htbp]
    \centering
    \caption{Ablación del número de pasos de puntuación~$T$ del DDPM en
             SICAP-C1 y SICAP-C12. Se reporta el AUROC medio para cada
             combinación de \textit{score} y~$T$. Un valor pequeño de~$T$
             reduce el coste computacional con pérdida marginal de rendimiento.}
    \label{tab:ablation_t}
    \input{tables/table_ablation_t_sicap.tex}
\end{table}

% ─────────────────────────────────────────────────────────────────────────────
\section{Análisis de errores y limitaciones}
\label{sec:res-limitaciones}
% ─────────────────────────────────────────────────────────────────────────────

Los resultados obtenidos permiten identificar tres fuentes principales de limitación
que condicionan la generalización de las conclusiones de este trabajo.

\paragraph{Acoplamiento entre representación y método.}~\\
Para SICAP, todos los métodos —incluyendo los modelos generativos— operan sobre
las características UNI y no sobre los píxeles originales. Esto significa que los
modelos generativos aprenden la distribución de una representación prefijada por un
modelo externo, y no la distribución de los datos crudos. Si el modelo UNI comprime
o elimina información relevante para la distinción entre Gleason~1 y Gleason~4,
ese límite superior de discriminabilidad afecta por igual a todos los métodos
comparados. La ablación propuesta no permite, por tanto, concluir si los modelos
generativos serían más o menos ventajosos en un escenario de entrenamiento
extremo a extremo sobre las imágenes originales.

\paragraph{Definición de la distribución ID.}~\\
Las configuraciones SICAP-C1 y SICAP-C12 definen la normalidad mediante etiquetas
de grado Gleason, que son una anotación patológica experta con variabilidad
inter-observador conocida~\cite{allsbrook_interobserver_2001}. Un modelo entrenado
con SICAP-C1 aprende lo que una anotación determinada considera Gleason~1; las
muestras de Gleason~2 anotadas por un patólogo diferente podrían ser
estadísticamente similares a las muestras ID de entrenamiento, lo que generaría
\textit{scores} bajos para una OOD que debería puntuar alto. Esta fuente de
ambigüedad es inherente a la definición del problema y no puede resolverse mediante
ningún método de modelado.

\paragraph{Coste computacional de evaluación y escalabilidad.}~\\
Los modelos DDPM en modo \textit{multi-step} requieren $T$ pasadas por el
\textit{denoiser} para puntuar cada muestra. Para $T{=}10$ y modelos MLP de
dimensión moderada, el coste es tolerable; sin embargo, el mismo protocolo aplicado
sobre modelos UNet de mayor capacidad o sobre conjuntos de evaluación grandes
podría resultar prohibitivo en escenarios de despliegue real. Los modos basados en
el residuo de predicción de ruido (Mahalanobis y k-NN) requieren, además, el
almacenamiento y la inversión de la matriz de covarianza de los residuos del
entrenamiento, cuyo coste escala cuadráticamente con la dimensión de entrada.
Para dimensiones superiores a las de SICAP~($d{=}1\,024$), se recomienda explorar
aproximaciones de baja dimensión como la descomposición en componentes principales
del espacio de residuos.

\paragraph{Variabilidad entre semillas.}~\\
Las barras de error de la figura~\ref{fig:sicap_methods_comparison} revelan que
la variabilidad del AUROC entre semillas no es despreciable, especialmente para
los modelos DDPM en SICAP-C12. Este fenómeno puede atribuirse a la sensibilidad
del proceso de difusión al estado inicial del generador de números aleatorios durante
la construcción de la referencia OOD, donde el ruido fijo $\boldsymbol{\varepsilon}$
semillado condiciona los estadísticos de normalización de los \textit{scores}
\textit{multi-step}. En futuros trabajos, promediar los \textit{scores} sobre
múltiples realizaciones de ruido podría reducir esta variabilidad a coste de un
incremento lineal en el tiempo de evaluación.

Los \textit{scores} OOD de MNIST y PathMNIST, las galerías extendidas de
reconstrucción y la tabla completa de SICAP se recogen en el
apéndice~\ref{ch:apendice-figuras}.

\endinput
